\chapter{مقدمه و پیش‌نیازها}
\thispagestyle{empty}

\section{مقدمه}
با گسترش روزافزون هوش مصنوعی و به ویژه مدل‌های زبانی بزرگ \lr{(LLMs)}، چشم‌انداز سیستم‌های تعاملی و دستیارهای هوشمند در حوزه‌های مختلف دگرگون شده است. یکی از حوزه‌های بسیار حساس و نیازمند دقت بالا، حوزه پزشکی و به طور خاص تخصص اورولوژی است. دسترسی سریع، دقیق و مستند به اطلاعات تخصصی می‌تواند هم برای کادر درمان به عنوان یک ابزار کمک‌تصمیم‌گیری و هم برای بیماران به منظور افزایش آگاهی، بسیار ارزشمند باشد. با این حال، استفاده مستقیم از هوش مصنوعی در این حوزه نیازمند غلبه بر چالش‌هایی نظیر تولید اطلاعات نادرست یا غیرمستند است. این رساله به دنبال طراحی یک سامانه پرسش و پاسخ دوزبانه است که بتواند با بالاترین دقت به نیازهای اطلاعاتی در این حوزه پاسخ دهد.

\section{مفاهیم پایه}
پردازش زبان طبیعی \lr{(NLP)} شاخه‌ای از هوش مصنوعی است که به تعامل بین رایانه و زبان انسانی می‌پردازد. در سال‌های اخیر، مدل‌های زبانی بزرگ توانسته‌اند درک عمیقی از الگوهای زبانی به دست آورند و متونی روان و مشابه انسان تولید کنند. با این وجود، یکی از مشکلات بنیادین این مدل‌ها پدیده‌ای به نام توهم \lr{(Hallucination)} است. توهم زمانی رخ می‌دهد که مدل زبانی پاسخی تولید می‌کند که از نظر نگارشی کاملاً صحیح و قانع‌کننده به نظر می‌رسد، اما از نظر علمی و واقعی نادرست است. در محیط‌های حساس پزشکی، این پدیده می‌تواند خطرات جدی به همراه داشته باشد؛ بنابراین نیازمند سازوکاری هستیم که مدل را ملزم کند تنها بر اساس مستندات معتبر پاسخ دهد.

\section{\lr{RAG}}
برای حل مشکل توهم و همچنین به‌روزرسانی دانش مدل‌های زبانی بدون نیاز به آموزش مجدد و پرهزینه، معماری تولید افزوده با بازیابی یا \lr{RAG} معرفی شده است. در این رویکرد، فرآیند پاسخ‌گویی به دو مرحله اصلی تقسیم می‌شود: ابتدا سیستم در یک پایگاه دانش محلی و معتبر (مانند کتب مرجع پزشکی) جستجو کرده و مرتبط‌ترین بخش‌های متن را بازیابی می‌کند. سپس، این متون بازیابی‌شده به عنوان زمینه \lr{(Context)} در اختیار مدل زبانی قرار می‌گیرند تا پاسخ نهایی را دقیقاً بر اساس آن‌ها استنتاج و تولید کند. این روش نه تنها دقت و صحت پاسخ‌ها را به شدت افزایش می‌دهد، بلکه امکان ارائه ارجاع دقیق به منابع اصلی را نیز فراهم می‌سازد.
